\documentclass[11pt,a4paper]{article}

\usepackage[margin=2.3cm]{geometry}
\usepackage[T1]{fontenc}
\usepackage{mathptmx}
\usepackage{amsmath,amssymb}
\usepackage{booktabs,tabularx}
\usepackage{xcolor}
\usepackage[shortlabels]{enumitem}
\usepackage{listings}
\usepackage{microtype}
\usepackage[hidelinks]{hyperref}

\setlength{\parskip}{0.4em}
\setlength{\parindent}{0pt}
\setlist{itemsep=0.15em,topsep=0.25em}
\lstset{basicstyle=\ttfamily\footnotesize, columns=fullflexible, frame=single, framerule=0.3pt,
  rulecolor=\color{black!40}, breaklines=true, keepspaces=true, xleftmargin=2pt, xrightmargin=2pt}
\newcommand{\file}[1]{\texttt{#1}}

\title{\textbf{Handoff: finishing the QPASE project report}}
\author{From Iman Kalyan Chakraborty (Person A) to Person B}
\date{\today}

\begin{document}
\maketitle
\raggedright

\section{Your job}

We are writing the project report together, in the style of the two sample reports in \file{sample\_reports/}.
I wrote the first half: the QPASE reproduction. You write the second half: our counter-based search index. When
you finish, the report is complete. The report is \file{report/main.tex}, and the compiled PDF is
\file{report/build/main.pdf}.

\begin{center}
\small
\begin{tabularx}{\linewidth}{lXll}
\toprule
\textbf{Section} & \textbf{Content} & \textbf{Owner} & \textbf{State} \\
\midrule
Abstract & summary of both halves & Person A, you revise & draft \\
1 Introduction & background, problem, contributions & Person A & done \\
2 Literature review and preliminaries & related work, notation, building blocks & Person A & done \\
3 The QPASE scheme & the paper's protocol, phase by phase & Person A & done \\
4 Implementation & modules, ring arithmetic, gaps in the paper & Person A & done \\
5 Evaluation of the reproduced scheme & costs, traffic, two security findings & Person A & done \\
6 Proposed scheme & the counter-based index & \textbf{you} & stub \\
7 Analysis of the proposed scheme & correctness, leakage, forward privacy, attacks & \textbf{you} & stub \\
8 Evaluation of the proposed scheme & notebook 07 results, search scaling & \textbf{you} & stub \\
9 Comparative study & both schemes; related password-based schemes & \textbf{you} & stub \\
10 Conclusion and future directions & summary and open work & \textbf{you} & stub \\
Appendix A & how every number is reproduced & Person A & done \\
\bottomrule
\end{tabularx}
\end{center}

Each of your sections already exists in \file{report/sections/} with its final heading, its label and a red
``To do'' box. Delete the box when you write the section. Sections 1 to 5 already cross-reference your labels
(\file{sec:counter-index}, \file{sec:counter-analysis}, \file{sec:counter-evaluation}, \file{sec:comparison},
\file{sec:conclusion}), so keep those labels.

\section{Set up and build}

You need the project environment and a LaTeX installation with \file{latexmk}, \file{pdflatex} and \file{bibtex}.
The report uses only standard packages, with no \file{biblatex}, \file{siunitx} or algorithm packages.
\begin{enumerate}
  \item In the repository root, run \file{uv sync} if \file{.venv/} does not exist yet.
  \item To build the report, run this in \file{report/}:
\begin{lstlisting}
latexmk main.tex
\end{lstlisting}
        The PDF is \file{report/build/main.pdf}. \file{report/.latexmkrc} sets PDF mode and the \file{build/}
        directory, which git ignores.
  \item To check everything at once, run this in the repository root:
\begin{lstlisting}
.venv/bin/python report/tools/check_report.py
\end{lstlisting}
        It checks that the generated numbers match the notebooks, builds the report, reports any error,
        undefined reference or undefined citation, and lists the open ``To do'' boxes.
\end{enumerate}
Do not rerun the notebooks. The report quotes the outputs committed in \file{notebooks/}, which all come from one
run on one machine. A rerun changes every timing, and the README and the report would then need refreshing too.

\section{What to write}

Target about 8 to 10 pages for Sections 6 to 10 together. The report is 20 pages now, stubs included, so it would
end near 28 to 30 pages, longer than both samples (9 and 21 pages). Your primary sources are:
\begin{itemize}
  \item \file{docs/counter-index-spec.md}: the full design, including the decisions the team made while building it.
  \item Slides 14 to 20 of \file{QPASE\_Paper\_Presentation\_counter\_index-2.pdf}: the motivation and the worked
        example.
  \item The code: \file{src/ris\_sem\_project/counter\_index.py}, \file{parties/counter\_index\_server.py},
        \file{parties/counter\_index\_user.py} and \file{protocol.py} (\file{ci\_login}, \file{ci\_outsource},
        \file{ci\_retrieve}).
  \item The README sections ``Counter-based search index'' and ``Results'', and notebooks 07 and 08.
\end{itemize}

\subsection{Section 6: Proposed scheme}

Write it like Section 3, so a reader can compare the two schemes phase by phase.
\begin{enumerate}
  \item \textbf{Motivation.} The two problems of QPASE's search: the server keeps $\mathit{dsk}$ and can test every
        later upload with it (no forward privacy), and a search tries $\mathit{dsk}$ on every stored item.
  \item \textbf{Design.} The keys $K_w = \mathrm{HKDF}(K_u, w)$, $k_{ptr}$ and $k_{ctr}$, none of which leaves the
        device. The label of the $j$-th file under $w$, the encrypted entry that holds the item id and the file's
        hash, and the counter table: sealed under $k_{ctr}$, padded to 1024-byte blocks, versioned, signed and kept
        on all $N$ servers. A small notation table helps, like \texttt{tab:notation} in Section 2.
  \item \textbf{Phases.} Login step L4, the new Outsource and the new Retrieve (R1 to R5), each in a
        \texttt{protocol} box (snippet in Section~\ref{sec:snippets}). Register, L1 to L3 and the key update do not change,
        so say so in one sentence.
  \item \textbf{Related designs.} The index is a counter-based encrypted dictionary in the style of Cash et
        al.~\verb|\cite{cash2014dynamic}|. Forward privacy as a goal comes from Stefanov et
        al.~\verb|\cite{stefanov2014practical}| and Bost~\verb|\cite{bost2016sophos}|.
\end{enumerate}

\subsection{Section 7: Analysis of the proposed scheme}

\begin{enumerate}
  \item \textbf{Correctness.} Why a search returns exactly the files under $w$: labels $1$ to $c_w$, the label
        inside each entry, the hash check on each fetched file, and duplicates kept once.
  \item \textbf{Leakage.} The table ``What the server learns'' from the spec or the README, with a sentence per row.
  \item \textbf{Forward privacy.} The server never holds $K_w$, so it cannot compute the label of a later upload, and
        a stored search request finds only older files. Note that the scheme does not give backward privacy and
        that the codebase has no deletion.
  \item \textbf{Attacks.} One paragraph each, as Section 4.3 of sample 1 does: a server moving an entry to another
        label, swapping a stored file, answering with another keyword's row (R3 drops labels it did not ask for),
        serving a stale or forged counter table (signature and version), and replaying a search request.
  \item \textbf{Costs we accept.} From slide 20 and the spec: the counter table grows with the number of keywords,
        Login gains a download, Retrieve takes two round trips, two devices at once must sync the counter, and the
        server still sees which files a search fetches.
  \item \textbf{Quantum resistance.} The index adds no new primitive. It uses the lattice PRF and HKDF, ML-DSA-44
        and AES-256-GCM that QPASE already uses. Argue this informally; we have no proof.
\end{enumerate}

\subsection{Section 8: Evaluation of the proposed scheme}

\begin{enumerate}
  \item \textbf{Functional results from notebook 07}, in the style of Section 5.2: sizes (\verb|\CiLabelBytes|,
        \verb|\CiEntryBytes|, \verb|\CiOutsourceRequestBytes|, \verb|\CiSealedBytes|, \verb|\CiUploadBytes|,
        \verb|\CiSearchRequestBytes|, \verb|\CiFetchRequestBytes|), the lookups per server (\verb|\CiLookups|),
        the forward-privacy demonstration (a replayed old request finds \verb|\CiReplayItems|, a new search finds
        \verb|\CiNewSearchItems|), and a new device getting the counter table (\verb|\CiNewDeviceVersions|).
  \item \textbf{Search time against storage size from notebook 08.} The figure \file{search\_scaling.png} and the
        generated table \file{generated/tables/search\_scaling.tex} are ready (snippets in Section~\ref{sec:snippets}).
        Headline values: the counter index's slowest server took \verb|\ScalingIndexMin| to
        \verb|\ScalingIndexMax|, the paper's \verb|\ScalingPaperMin| to \verb|\ScalingPaperMax|, a speed-up of
        \verb|\ScalingSpeedupMin| to \verb|\ScalingSpeedupMax| (median \verb|\ScalingSpeedupMedian|), with
        \verb|\ScalingFailures| failures over \verb|\ScalingChecks| checked runs.
  \item \textbf{Explain the result honestly.} The speed-up does not come mainly from avoiding a slow scan: a
        failed AES-GCM check on a non-matching item is cheap. It comes from the paper's server running one lattice
        PRF chain (\verb|\AxMedian|) per matching item, while the counter index's server only looks labels up.
        The scan still shows. In the generated table the paper's slowest server rises at $10^5$ files at every
        $N$ (for example 537 to 721~ms at $N = 3$), while the counter index stays flat. Describe that rise from the
        table; do not call the paper's search time constant. The user's time is about the same in both schemes.
\end{enumerate}

\subsection{Section 9: Comparative study}

\begin{enumerate}
  \item \textbf{QPASE against the counter index.} A table with one row per property: whether the search key
        leaves the device, server work per search, forward privacy, what a server learns, round trips per search,
        storage per file, and the user's cost per search. Use the macros above for the numbers.
  \item \textbf{QPASE against related password-based schemes.} Reproduce the paper's Table V (page 14 of the
        paper PDF) as a table, citing the paper. My transcription is below. Check it against the paper before you
        use it. Every scheme in it has an entry in \file{references.bib}; the keys are in Section~\ref{sec:citations}.
\end{enumerate}

{\small
\begin{tabularx}{\linewidth}{Xllllllll}
\toprule
\textbf{Scheme} & \textbf{Threshold} & \textbf{Passwords} & \textbf{Model} & \textbf{Tech.} & \textbf{Quantum} &
\textbf{Retrieval} & \textbf{Rounds} & \textbf{Server / user} \\
\midrule
Bagherzandi et al.\ CCS'11 & $(t,N)$ & uniform & ROM & HE & no & no & 2 & $16 / 33\,T_{exe}$ \\
Camenisch et al.\ CCS'12 & $(2,2)$ & n/a & UC & HE & no & no & 1 & $26 / 19\,T_{exe}$ \\
Jarecki et al.\ ASIACRYPT'14 & $(t,N)$ & uniform & ROM & OPRF & no & no & 1 & $4 / 11\,T_{exe}$ \\
Yi et al.\ ESORICS'15 & $(t,N)$ & uniform & UC & HE & no & no & 1 & $12 / 7\,T_{exe}$ \\
Jarecki et al.\ EuroS\&P'16 & $(t,N)$ & uniform & UC & OPRF & no & no & 1 & $1 / 4\,T_{exe}$ \\
Jarecki et al.\ ACNS'17 & $(t,N)$ & uniform & UC & OPRF & no & no & 1 & $1 / 2\,T_{exe}$ \\
Das et al.\ AsiaCCS & $(N,N)$ & uniform & UC & OPRF & no & no & 2 & $4 / 10\,T_{exe}$ \\
Chen et al.\ IJIS'21 (PASE) & $(2,2)$ & uniform & ROM & HE & no & yes & 1 & $8 / 3\,T_{exe}$ \\
Roy et al.\ ACNS'21 & $(t,N)$ & uniform & ROM & FHE & yes & no & 1 & $0 / 0$ \\
Jiang et al.\ TSC'23 (QPause) & $(t,N)$ & Zipf & ROM & FHE & yes & no & 1 & $0 / 0$ \\
QPASE & $(t,N)$ & Zipf & ROM & OPRF & yes & yes & 1 & $0 / 0$ \\
\bottomrule
\end{tabularx}
}

$T_{exe}$ is one exponentiation. HE is homomorphic encryption and FHE fully homomorphic encryption. The paper
lists Das et al.\ as ``ASIACCS'20'' in the table, while its reference list dates the paper 2022; cite the reference
list's version.

\subsection{Section 10: Conclusion and future research directions}

Summarise both halves in a few paragraphs: QPASE works as specified once the gaps are filled; the PRF sets its
cost; PARAM~II in this ring is unsafe; the linear key-update commitments reveal the shares; the counter index gives forward
privacy and faster server-side search at the same user cost. Then list future work, for example: a hiding,
binding commitment for the key update; parameters for PARAM~II in a safe ring or in plain LWE; a faster
quantum-safe PRF for labels and HKDF; multi-keyword search (Section IV-G of the paper); backward privacy and
deletion; syncing the counter table across devices; constant-time sampling; and an implementation in a compiled
language.

\subsection{Abstract and author names}

When your sections are written, revise the abstract in \file{sections/00-abstract.tex} so it matches them, and
delete its ``To do'' box. In \file{main.tex}, replace \verb|[Partner name]| and \verb|[Roll number]| with yours.

\section{Numbers, figures and citations}

\subsection{Never type a measured value}

Every measured value is a LaTeX macro in \file{report/generated/numbers.tex}, generated from the saved notebook
outputs by \file{report/tools/make\_numbers.py}. Each line's comment names the notebook cell it came from, and
Appendix A of the report lists them all. Write \verb|\CiSealedBytes|, never \verb|1072 B|. This keeps the report
in step with the notebooks and the README, and it is how we meet the project rule that every number can be traced
to a notebook cell.

If you need a value that has no macro yet, add one entry to the \texttt{NUMBERS} list in
\file{make\_numbers.py}: a macro name (letters only), the notebook number, the cell number, and a regular
expression whose first group captures the value exactly as the cell prints it. For example:
\begin{lstlisting}
Number("CiSearchRequestBytes", "07", 9, r"LabelSearchRequest, (\d+ B)", "label search request for c_w = 3"),
\end{lstlisting}
Then run \file{.venv/bin/python report/tools/make\_numbers.py}. It stops with an error if the pattern does not
match, so a wrong pattern cannot slip into the report. The paper's own numbers, such as its Table IV, are not
generated: write them with a citation, as Section 5.4 does.

\subsection{Snippets}
\label{sec:snippets}

A protocol phase, as in Section 3:
\begin{lstlisting}
\begin{protocol}{Retrieve with the counter index}
\begin{enumerate}[R1.]
  \item If $c_w = 0$, return no files. Otherwise ...
  \item ...
\end{enumerate}
\end{protocol}
\end{lstlisting}

The search-scaling figure and table. Figures are found in \file{docs/figures/}:
\begin{lstlisting}
\begin{figure}[ht]
\centering
\includegraphics[width=\linewidth]{search_scaling.png}
\caption{Median search time at the slowest server ... Drawn by
\texttt{notebooks/08\_search\_scaling.ipynb}, cell 7.}
\label{fig:search-scaling}
\end{figure}

\begin{table}[ht]
\centering
\small
\caption{... Source: \texttt{notebooks/08\_search\_scaling.ipynb}, cell 7.}
\label{tab:search-scaling}
\resizebox{\linewidth}{!}{\input{generated/tables/search_scaling}}
\end{table}
\end{lstlisting}

Cross-references use \verb|\cref{...}|, which prints ``Section 6'' or ``Table 9'' for you.

\subsection{Citations}
\label{sec:citations}

\file{report/references.bib} holds every work the report needs. Most entries copy the QPASE paper's own
reference list; a comment above each entry says so. The keys you are most likely to use are
\verb|jiang2024qpase|, \verb|chen2021pase|, \verb|curtmola2006sse|, \verb|cash2014dynamic|,
\verb|stefanov2014practical|, \verb|bost2016sophos|, \verb|bagherzandi2011pps|, \verb|camenisch2012pass|,
\verb|jarecki2014round|, \verb|yi2015tpass|, \verb|jarecki2016highly|, \verb|jarecki2017toppss|,
\verb|das2022dpase|, \verb|roy2021pps| and \verb|jiang2024qpause|. Four entries are not from the paper's list
and are marked so: \verb|lyubashevsky2010ring|, \verb|fips204|, \verb|dworkin2007gcm| and
\verb|bost2016sophos|. Check those four against the publishers' pages before submission. Never add a citation
from memory without checking it.

\section{Writing rules}

Match Sections 1 to 5 so the two halves read as one report.
\begin{itemize}
  \item Write in the first person plural, present tense, active voice: ``The server looks up each label.''
  \item One idea per sentence. Short sentences.
  \item No em dashes. Use a full stop or a comma instead.
  \item Name the real thing: the macro, the file, the notebook cell, the protocol step.
  \item Every claim carries its evidence: a macro, a cited table or figure, or a citation. Mark an inference as
        one (``we infer \dots''), as Section 5.4 does.
  \item Use the paper's notation (Table 1 of the report) and the spec's names ($K_w$, $k_{ptr}$, $k_{ctr}$, $c_w$).
\end{itemize}

\section{When you are done}

The report is finished when all of these hold:
\begin{enumerate}
  \item This command exits with status 0 (it fails while any ``To do'' box or name placeholder remains):
\begin{lstlisting}
.venv/bin/python report/tools/check_report.py --final
\end{lstlisting}
  \item The abstract matches Sections 6 to 10.
  \item You have checked the four citations that do not come from the paper's reference list.
  \item You have read the whole PDF once, front to back, as a reviewer would.
\end{enumerate}
Then send \file{report/build/main.pdf}.

\end{document}
